\section{Related Work}
\label{sec:related}

\paragraph{Apportionment.}
The Huntington-Hill method \citep{huntington1921methods, hill1911method}
allocates congressional seats using a geometric-mean priority rule: a state
with $n$ seats receives the next seat before any state with fewer if its
priority $P / \sqrt{n(n+1)}$ (where $P$ is population) exceeds all competitors.
\citet{balinski2001fair} provide the canonical treatment of apportionment
methods and their fairness properties. Crucially, no prior work connects the
arithmetic structure of the resulting seat counts to the geometry of the
districts those seats represent.

\paragraph{Graph partitioning.}
The $k$-way graph partitioning problem --- partition the vertices of a weighted
graph into $k$ balanced parts minimising edge cut --- is NP-hard
\citep{garey1979computers}. The METIS multilevel scheme
\citep{karypis1998fast} provides practical heuristics for both recursive
bisection ($k = 2^d$) and direct $k$-way partitioning. Our prior work
\citep{dellaLibera2026mec} established minimum-edge-cut (MEC) redistricting
using recursive bisection and characterised its seed sensitivity.
PFR generalises MEC by using \emph{prime}-way splits, not only binary ones.

\paragraph{Hierarchical redistricting.}
\citet{cohen2021making} and \citet{hegeman2020recursive} study recursive
district structures but do not connect the recursion depth or branching factor
to the prime factorization of the seat count. Existing hierarchical methods
treat the bisection tree structure as an implementation detail, not as a
mathematically-determined property of the seat count.

\paragraph{Constitutional basis.}
\citet{rucho2019} (Rucho v.\ Common Cause) held that partisan gerrymandering
claims are non-justiciable in federal court but did not foreclose congressional
regulation of the manner of elections. \citet{moore2023} (Moore v.\ Harper)
rejected the independent state legislature doctrine, confirming that both state
and federal law constrain redistricting. PFR's constitutional argument anchors
redistricting in Article~I~\S2 (apportionment) rather than \S4 (Manner),
a distinction unexplored in the legal literature.

\paragraph{Algorithmic redistricting.}
\citet{duchin2019gerrychain} propose Markov-chain sampling of the redistricting
plan space. \citet{herschlag2020quantifying} use sampling to quantify
gerrymandering in North Carolina. \citet{chikina2017assessing} develop
statistical significance tests for partisan plans. PFR differs from all of
these by not sampling a plan space at all: the plan is uniquely determined by
arithmetic and geography.
